\subsection{Обоснование редуцированного представления состояния}

Используемое в работе редуцированное представление состояния не претендует на полное описание внутренней структуры HANK-модели. Его задача состоит в том, чтобы выделить те скрытые компоненты, которые одновременно имеют содержательный экономический смысл, поддаются оцениванию по наблюдаемым данным и значимы для выбора процентной ставки.

Обоснование такого представления опирается на три критерия. Во-первых, каждый компонент редуцированного состояния сопоставляется с одним из каналов денежно-кредитной трансмиссии в HANK-среде. Во-вторых, состояние проверяется на прогностическую достаточность: оно должно содержать информацию о будущей инфляции, разрыве выпуска и будущей функции потерь. При этом в прогнозной валидации используется только информационный набор, доступный до выбора текущей ставки, включая лагированную процентную ставку; текущая выбранная ставка исключается во избежание информационной утечки. В-третьих, состояние проверяется на управленческую достаточность: ранжирование правил, полученное в редуцированной постановке, должно сохраняться при пропуске соответствующих траекторий ставки через HANK-проекцию.

\begin{table}[htbp]
\centering
\small
\caption{Экономическая интерпретация компонент редуцированного состояния}
\label{tab:reduced_state_components}
\begin{tabular}{p{0.22\linewidth}p{0.34\linewidth}p{0.34\linewidth}}
\toprule
Компонент & Экономический смысл & Канал HANK \\
\midrule
Разрыв естественной ставки & Скрытая компонента межвременного условия и нейтральной ставки & Межвременной канал денежно-кредитной трансмиссии \\
Производственная компонента & Скрытый сдвиг производительности и трудового дохода & Доходный и производственный канал \\
Фискальная компонента & Скрытый сдвиг фискальной позиции и трансфертов & Фискальный канал и доходы домохозяйств \\
Разрыв инфляции & Отклонение инфляции от стационарного уровня & Номинальные жесткости и стандартный стабилизационный мотив \\
Разрыв выпуска & Отклонение выпуска от стационарного уровня & Совокупный спрос, занятость и трудовой доход \\
Доля низколиквидных домохозяйств & Скрытый индикатор напряженности ликвидных балансов & Распределительный и ликвидностный канал \\
Средняя предельная склонность к потреблению & Скрытый индикатор чувствительности потребления к доходу & Канал гетерогенных предельных склонностей к потреблению \\
Вероятность стрессового режима & Апостериорная вероятность скрытого режима с иной трансмиссией & Режимная неопределенность и изменение силы передачи политики \\
\bottomrule
\end{tabular}
\end{table}

\begin{table}[htbp]
\centering
\small
\caption{Прогностическая достаточность альтернативных представлений состояния}
\label{tab:forecast_sufficiency}
\begin{tabular}{p{0.30\linewidth}rrrr}
\toprule
Представление состояния & Инфляция & Разрыв выпуска & Потери & Число сценариев \\
\midrule
Только наблюдаемые переменные & 0.781 & 0.744 & 0.057 & 4 \\
Макроэкономическое состояние & 0.784 & 0.750 & 0.050 & 4 \\
Макроэкономическое состояние + скрытый режим & 0.784 & 0.750 & 0.182 & 4 \\
Распределительно расширенное состояние & 0.783 & 0.746 & 0.180 & 4 \\
\bottomrule
\end{tabular}
\end{table}

\begin{table}[htbp]
\centering
\small
\caption{Сценарная прогностическая достаточность для будущей функции потерь}
\label{tab:future_loss_scenario_forecast}
\begin{tabular}{p{0.32\linewidth}rrrr}
\toprule
Сценарий & Наблюдаемые переменные & Макроэкономическое состояние & Макроэкономическое состояние + скрытый режим & Распределительно расширенное состояние \\
\midrule
Базовый макроэкономический набор × умеренная разделимость режимов & 0.056 & 0.055 & 0.044 & 0.037 \\
Базовый макроэкономический набор × высокая разделимость режимов & 0.060 & 0.028 & 0.060 & 0.059 \\
Ограниченный информационный набор × умеренная разделимость режимов & 0.058 & 0.068 & 0.336 & 0.339 \\
Ограниченный информационный набор × высокая разделимость режимов & 0.053 & 0.048 & 0.288 & 0.287 \\
\bottomrule
\end{tabular}
\end{table}

\begin{table}[htbp]
\centering
\small
\caption{Сохранение ранжирования правил при проверке через HANK-проекцию}
\label{tab:policy_ranking_validation}
\begin{tabular}{p{0.30\linewidth}p{0.25\linewidth}rrrr}
\toprule
Сценарий & Правило & Ранг в редуцированной среде & Ранг HANK & Масштаб & Потери HANK \\
\midrule
Базовый макроэкономический набор × умеренная разделимость режимов & Историческое по наблюдениям & 1.0 & 1.0 & 1.00 & 1.267e-06 \\
Базовый макроэкономический набор × умеренная разделимость режимов & Оптимизированное линейное & 2.0 & 2.0 & 1.00 & 1.097e-05 \\
Базовый макроэкономический набор × умеренная разделимость режимов & Классическое правило & 3.0 & 3.0 & 0.25 & 1.248e-04 \\
Базовый макроэкономический набор × высокая разделимость режимов & Историческое по наблюдениям & 1.0 & 1.0 & 1.00 & 1.924e-06 \\
Базовый макроэкономический набор × высокая разделимость режимов & Оптимизированное линейное & 2.0 & 2.0 & 1.00 & 1.453e-05 \\
Базовый макроэкономический набор × высокая разделимость режимов & Классическое правило & 3.0 & 3.0 & 0.25 & 1.873e-04 \\
Ограниченный информационный набор × умеренная разделимость режимов & Историческое по наблюдениям & 1.0 & 1.0 & 1.00 & 1.267e-06 \\
Ограниченный информационный набор × умеренная разделимость режимов & Оптимизированное линейное & 2.0 & 2.0 & 1.00 & 9.935e-06 \\
Ограниченный информационный набор × умеренная разделимость режимов & Классическое правило & 3.0 & 3.0 & 0.25 & 1.346e-04 \\
Ограниченный информационный набор × высокая разделимость режимов & Историческое по наблюдениям & 1.0 & 1.0 & 1.00 & 1.924e-06 \\
Ограниченный информационный набор × высокая разделимость режимов & Оптимизированное линейное & 2.0 & 2.0 & 1.00 & 1.240e-05 \\
Ограниченный информационный набор × высокая разделимость режимов & Классическое правило & 3.0 & 3.0 & 0.25 & 1.750e-04 \\
\bottomrule
\end{tabular}
\end{table}

\begin{table}[htbp]
\centering
\small
\caption{Попарные разности потерь в common-scale HANK-проекции}
\label{tab:common_scale_pairwise_intervals}
\begin{tabular}{p{0.21\linewidth}p{0.22\linewidth}p{0.11\linewidth}rrrr}
\toprule
Сценарий & Сравнение & Число траекторий & $\Delta J^{red}$ & $\Delta J^{HANK}$ & 95\%-й ДИ & Нуль в ДИ \\
\midrule
Базовый макроэкономический набор × умеренная разделимость режимов & Классическое правило -- Историческое по наблюдениям & 3 & 1.47e-03 & 4.03e-04 & [3.13e-04; 5.23e-04] & Нет \\
Базовый макроэкономический набор × умеренная разделимость режимов & Классическое правило -- Оптимизированное линейное & 3 & 1.41e-03 & 3.97e-04 & [3.08e-04; 5.16e-04] & Нет \\
Базовый макроэкономический набор × умеренная разделимость режимов & Историческое по наблюдениям -- Оптимизированное линейное & 12 & -1.54e-04 & -1.85e-05 & [-2.80e-05; -1.08e-05] & Нет \\
Базовый макроэкономический набор × высокая разделимость режимов & Классическое правило -- Историческое по наблюдениям & 1 & 1.09e-03 & 3.74e-04 & [3.74e-04; 3.74e-04] & Нет \\
Базовый макроэкономический набор × высокая разделимость режимов & Классическое правило -- Оптимизированное линейное & 1 & 1.07e-03 & 3.69e-04 & [3.69e-04; 3.69e-04] & Нет \\
Базовый макроэкономический набор × высокая разделимость режимов & Историческое по наблюдениям -- Оптимизированное линейное & 12 & -1.12e-04 & -2.11e-05 & [-2.98e-05; -1.26e-05] & Нет \\
Ограниченный информационный набор × умеренная разделимость режимов & Классическое правило -- Историческое по наблюдениям & 3 & 1.14e-03 & 3.85e-04 & [3.21e-04; 4.57e-04] & Нет \\
Ограниченный информационный набор × умеренная разделимость режимов & Классическое правило -- Оптимизированное линейное & 3 & 1.09e-03 & 3.79e-04 & [3.15e-04; 4.51e-04] & Нет \\
Ограниченный информационный набор × умеренная разделимость режимов & Историческое по наблюдениям -- Оптимизированное линейное & 12 & -1.34e-04 & -1.69e-05 & [-2.59e-05; -1.01e-05] & Нет \\
Ограниченный информационный набор × высокая разделимость режимов & Классическое правило -- Историческое по наблюдениям & 3 & 9.35e-04 & 3.36e-04 & [2.06e-04; 4.77e-04] & Нет \\
Ограниченный информационный набор × высокая разделимость режимов & Классическое правило -- Оптимизированное линейное & 3 & 9.20e-04 & 3.32e-04 & [2.05e-04; 4.72e-04] & Нет \\
Ограниченный информационный набор × высокая разделимость режимов & Историческое по наблюдениям -- Оптимизированное линейное & 12 & -9.02e-05 & -1.86e-05 & [-2.59e-05; -1.16e-05] & Нет \\
\bottomrule
\end{tabular}
\end{table}

\begin{table}[htbp]
\centering
\small
\caption{Чувствительность ранжирования к масштабу HANK-проекции}
\label{tab:common_scale_sensitivity}
\begin{tabular}{p{0.30\linewidth}rrr}
\toprule
Сценарий & Масштаб & Попарные знаки & Ранговая корреляция \\
\midrule
Базовый макроэкономический набор × умеренная разделимость режимов & 0.10 & 3/3 & 1.00 \\
Базовый макроэкономический набор × умеренная разделимость режимов & 0.25 & 3/3 & 1.00 \\
Базовый макроэкономический набор × высокая разделимость режимов & 0.10 & 3/3 & 1.00 \\
Базовый макроэкономический набор × высокая разделимость режимов & 0.25 & 3/3 & 1.00 \\
Ограниченный информационный набор × умеренная разделимость режимов & 0.10 & 3/3 & 1.00 \\
Ограниченный информационный набор × умеренная разделимость режимов & 0.25 & 3/3 & 1.00 \\
Ограниченный информационный набор × высокая разделимость режимов & 0.10 & 3/3 & 1.00 \\
Ограниченный информационный набор × высокая разделимость режимов & 0.25 & 3/3 & 1.00 \\
\bottomrule
\end{tabular}
\end{table}

Результаты показывают, что базовым представлением состояния в среднем следует считать макроэкономическое состояние с учётом скрытого режима, тогда как распределительно расширенное состояние разумно трактовать как HANK-расширение, а не как универсально доминирующий прогнозный набор. Главный прирост прогнозного качества для будущей функции потерь связан с учётом скрытого режима, а не с механическим добавлением всех распределительных компонент. Ранжирование правил сохраняется при проверке через HANK-проекцию. При этом таблица \ref{tab:policy_ranking_validation} должна интерпретироваться с учётом масштаба траекторий: если поле масштаба меньше единицы, соответствующая траектория прошла через переходный решатель HANK только после уменьшения амплитуды. Поэтому данная проверка является проверкой согласованности и численной мягкости траекторий, а не полной переоптимизацией правил в исходной HANK-модели. В таблице \ref{tab:common_scale_pairwise_intervals} специально вынесено число сопоставимых траекторий: для пары двух неклассических правил доступны все 12 отложенных траекторий, тогда как для сравнений с классическим правилом из-за сходимости HANK-проекции остаются лишь 1--3 траектории в зависимости от сценария. Дополнительные interval-оценки на seed-level HANK-проекциях подтверждают устойчивое превосходство двух неклассических правил над классическим правилом в доступном множестве сходящихся траекторий. При этом различие между двумя лучшими правилами существенно меньше по масштабу и должно интерпретироваться осторожнее, чем их отрыв от классического правила.
